% IA 642 Defensive Security
% Week 05 Module 01 Lab 5.1: Mapping Meridian's Campaign in ATT&CK Navigator and STIX
% Group: Unais Ali (E02805019) and Hafiz Usama (E02574092)
% Submit as: IA642_Lab5.1_Ali.pdf (+ three JSON files)
\documentclass[11pt,letterpaper]{article}

\usepackage[margin=0.75in]{geometry}
\usepackage{array}
\usepackage{booktabs}
\usepackage{tabularx}
\usepackage{enumitem}
\usepackage{parskip}
\usepackage{ragged2e}
\usepackage{titlesec}
\usepackage{graphicx}
\usepackage{float}
\usepackage{placeins}
\usepackage{microtype}
\usepackage{xcolor}
\usepackage{colortbl}
\usepackage{listings}
\usepackage{caption}
\usepackage{needspace}
\usepackage[colorlinks=true,linkcolor=blue,citecolor=blue,urlcolor=blue]{hyperref}
\usepackage{xurl}

\hypersetup{pdfauthor={Unais Ali; Hafiz Usama}, pdftitle={IA 642 Week 05 Lab 5.1 ATT&CK Navigator and STIX}}

\graphicspath{{captures/}{./}}

\definecolor{tblhead}{HTML}{1F4E79}
\definecolor{tblrow}{HTML}{F0F4F8}
\definecolor{codebg}{HTML}{F7FAFC}
\definecolor{okbox}{HTML}{E6F4EA}
\setlength{\parindent}{0pt}
\renewcommand{\arraystretch}{1.22}
\raggedbottom
\setlist[itemize]{leftmargin=1.2em,itemsep=0.22em,topsep=0.3em,parsep=0pt}
\setlist[enumerate]{leftmargin=1.3em,itemsep=0.28em,topsep=0.3em,parsep=0pt}
\newcolumntype{Y}{>{\RaggedRight\arraybackslash}X}

\titleformat{name=\section,numberless}{\large\bfseries\color{tblhead}\filright}{}{0em}{}
\titlespacing*{name=\section,numberless}{0pt}{0.85em}{0.3em}
\titleformat{name=\subsection,numberless}{\normalsize\bfseries\color{tblhead}\filright}{}{0em}{}
\titlespacing*{name=\subsection,numberless}{0pt}{0.55em}{0.2em}

\newcommand{\labpart}[1]{%
  \par\vspace{0.35em}%
  \Needspace{8\baselineskip}%
  \section*{#1}%
  \nopagebreak
}
\newcommand{\labstep}[1]{%
  \Needspace{6\baselineskip}%
  \subsection*{#1}%
  \nopagebreak
}

\lstset{
  basicstyle=\ttfamily\scriptsize,
  backgroundcolor=\color{codebg},
  frame=single,
  rulecolor=\color{tblhead},
  breaklines=true,
  breakatwhitespace=false,
  columns=fullflexible,
  keepspaces=true,
  showstringspaces=false,
  aboveskip=0.35em,
  belowskip=0.35em,
  xleftmargin=0.15em,
  framexleftmargin=0.2em
}

\newcommand{\capturebox}[2]{%
  \par\vspace{0.25em}%
  \noindent\colorbox{okbox}{\parbox{\dimexpr\linewidth-2\fboxsep\relax}{%
    \textbf{CAPTURE \#{#1}.} #2}}%
  \par\vspace{0.25em}%
}

\newcommand{\capfig}[1]{%
  \par\noindent\centering
  \includegraphics[width=0.94\textwidth,height=0.38\textheight,keepaspectratio]{#1}%
  \par\vspace{0.25em}%
}

\begin{document}

\begin{center}
{\LARGE\bfseries IA 642 Defensive Security}\\[0.25em]
{\Large Week 05 Module 01 Lab 5.1}\\[0.2em]
{\large Mapping Meridian's Campaign in ATT\&CK Navigator and STIX}\\[0.35em]
{\small Cyber Threat Intelligence $\cdot$ MITRE ATT\&CK Enterprise v19 $\cdot$ STIX 2.1}\\[0.25em]
{\small Repository:
\href{https://github.com/unaisshazan/ia642-week05-lab51-cti}{github.com/unaisshazan/ia642-week05-lab51-cti}}\\[0.25em]
\end{center}

\begin{tabularx}{\textwidth}{@{}l X l l@{}}
\toprule
\textbf{Group members} & Unais Ali & \textbf{EID} & E02805019 \\
 & Hafiz Usama &  & E02574092 \\
\textbf{Course} & IA 642 Defensive Security & \textbf{Week} & 05, Lab 5.1 \\
\textbf{Tools} & Navigator 5.3.2, STIX 2.1 & \textbf{Date} & October 2026 \\
\bottomrule
\end{tabularx}

\vspace{0.55em}

\section*{Lab overview}

This lab maps Meridian's evidence cards (E1--E4) onto MITRE ATT\&CK Enterprise v19 in ATT\&CK Navigator, tests lecture beacon mappings against ATT\&CK wording, compares the Meridian layer to FIN7 (G0046), writes a three-object STIX~2.1 bundle, validates with the starter-kit checkers, and argues whether the intrusion looks opportunistic or targeted. Submitted files: this PDF, \texttt{meridian\_layer.json}, \texttt{meridian\_overlap\_layer.json}, and \texttt{meridian\_stix.json}.

\begin{tabularx}{\textwidth}{@{}l Y@{}}
\toprule
\textbf{Deliverable} & \textbf{Checker / note} \\
\midrule
\texttt{meridian\_layer.json} & \texttt{check\_layer.py} prints OK (7 techniques) \\
\texttt{meridian\_overlap\_layer.json} & \texttt{check\_layer.py --overlap} (overlap fraction 1 of 67) \\
\texttt{meridian\_stix.json} & \texttt{check\_stix.py} RESULT: ALL CHECKS PASSED \\
This report & Q1--Q6, CAPTURE \#1--\#4, FIN7 profile, Challenge \\
\bottomrule
\end{tabularx}

\labpart{Task 1. Setup}

Python 3.11.9 virtual environment in \texttt{lab51/}, with pinned packages from \texttt{requirements.txt} (\texttt{stix2-validator} 3.2.0, \texttt{stix2} 3.0.2). Running \texttt{check\_stix.py} on the skeleton produced the expected failures before editing \texttt{meridian\_stix.json}, confirming the checker works.

\labpart{Part A. Navigator layer (Tasks 2--6)}

\labstep{Q1. ATT\&CK version and T1027 tactic column}

The layer uses MITRE ATT\&CK Enterprise v19 (Navigator 5.3.2, layer format 4.5). In v19, T1027 Obfuscated Files or Information appears under the \textbf{Stealth} tactic column (TA0005). Week~05 materials note that v19 renamed the former Defense Evasion column to Stealth and added Defense Impairment (TA0112). T1027 was scored in Stealth to match the Navigator matrix.

\labstep{Anchored techniques (Tasks 2--3)}

\begin{tabularx}{\textwidth}{@{}l l c Y@{}}
\toprule
\textbf{Technique} & \textbf{Tactic} & \textbf{Score} & \textbf{Evidence basis} \\
\midrule
T1558.003 Kerberoasting & Credential Access & 95 & E2: nine 4769 RC4 (0x17) for one SPN \\
T1027 Obfuscated Files & Stealth & 45 & E3 stand-in structure (not svcupdate\_v2.bin) \\
\bottomrule
\end{tabularx}

\textbf{T1558.003 comment used.} E2: nine Event ID 4769 requests for one SPN (\texttt{MSSQLSvc/sqlapp01...}), RC4 (0x17), about 3.5 minutes, from an ordinary logged-in account; ticket cracked offline. Matches ATT\&CK Kerberoasting.

\textbf{T1027 comment used.} E3 stand-in shows MZ marker, MeridianLoader banner, decode-stub-shaped bytes (\texttt{90 90 90 31 C0...}), and embedded C2 host. Fits Obfuscated Files parent, but the artifact is a defanged stand-in, so confidence stays possible/contested.

\labstep{Q2. Parent vs sub-technique for T1027}

I scored the parent T1027, not a sub-technique such as T1027.002 Software Packing. The deciding E3 facts are the decode-stub-shaped byte sequence and embedded strings inside a non-valid PE stand-in. That shows obfuscated/encoded content without proving a specific packing method. I did not score higher because the file is an IR-defanged stand-in with no working code; citing \texttt{svcupdate\_v2.bin} would be wrong because that was a YARA test file I built, not adversary evidence.

\labstep{Q3. Task 4 beacon mappings (E1)}

\begin{tabularx}{\textwidth}{@{}l l c Y@{}}
\toprule
\textbf{Candidate} & \textbf{Decision} & \textbf{Score} & \textbf{Reason (ATT\&CK wording + E1)} \\
\midrule
T1071.001 Web Protocols & Include & 75 & Application-layer C2 over web ports; 143 TLS sessions to :443 \\
T1573.002 Asymmetric Crypto & Reject & 0 & Ordinary TLS is inherent protocol protection, not added crypto \\
T1029 Scheduled Transfer & Reject & 0 & Timer is C2 check-in, not scheduled data exfiltration \\
T1036.005 (optional) & Include & 40 & \texttt{svcupdate.bin} / MeridianLoader naming (E3) \\
\bottomrule
\end{tabularx}

\textbf{T1071.001.} ATT\&CK describes adversaries who ``communicate using application layer protocols to avoid detection/network filtering by blending in with existing traffic.'' E1 shows 143 outbound TLS sessions to \texttt{sync-update.meridian-relay.test} on TCP~443 with about an 8-minute cadence. Without a packet capture, the HTTP (or other) payload inside TLS cannot be confirmed, so the score is strong but not 90+.

\textbf{T1573.002.} ATT\&CK requires asymmetric encryption ``to conceal command and control traffic rather than relying on any inherent protections provided by a communication protocol.'' E1 is ordinary TLS~1.2 with a publicly trusted certificate. That is inherent protocol protection. Parent T1573 is not a better fit for the same reason.

\textbf{T1029.} ATT\&CK Scheduled Transfer is about scheduling \emph{data exfiltration}. E1's timer shows C2 check-ins (about 540--610~B send / 1.2--1.4~KB receive), not evidence that Meridian data left on a schedule.

\labstep{Q4. Task 5: T1110.001 Password Guessing (E4)}

No. I scored T1110.001 at 0 for this campaign. Two deciding E4 facts: (1)~alerts are on \texttt{vpn-gw-01} and \texttt{corp-backup-01}, not BILL-WS-14 or MERIDIAN-DC01; (2)~no time, account, or address link to E1--E3, and \texttt{corp-backup-01}'s success matches a known nightly password-change re-auth. Evidence that would change the answer: the same source IP or username appearing in E2's 4769 events, or a successful login on BILL-WS-14 / DC01 tied to the VPN failures.

\labstep{CAPTURE \#1: finished layer and check\_layer.py}

\capturebox{1}{Finished Meridian layer in ATT\&CK Navigator (show IDs on), plus full \texttt{check\_layer.py} output ending in OK.}

\capfig{capture1_navigator_layer.png}

\captionof{figure}{CAPTURE \#1: Finished Meridian layer in ATT\&CK Navigator 5.3.2 (show IDs on). T1558.003 Kerberoasting scored high (green). Companion board below lists every annotated technique.}

\capfig{capture1_layer.png}

\captionof{figure}{CAPTURE \#1 companion: full technique ID, name, tactic, and score list for all seven annotated techniques.}

\begin{lstlisting}
Layer: 'Meridian Lab 5.1 - Ali & Usama'
domain: enterprise-attack   ATT&CK v19   Navigator 5.3.2   layer format 4.5

Technique   Score  Comment  Tactic(s)
T1027          45  yes      stealth
T1029           0  yes      exfiltration
T1036.005      40  yes      stealth
T1071.001      75  yes      command-and-control
T1110.001       0  yes      credential-access
T1558.003      95  yes      credential-access
T1573.002       0  yes      command-and-control

7 technique(s) annotated.
OK -- every technique has a score and a comment.
\end{lstlisting}

\labpart{Part B. FIN7 comparison (Task 7)}

I compared Meridian against FIN7 (G0046). When FIN7 was selected in Navigator, the technique-count badge showed \textbf{67} (ATT\&CK v19). Overlap expression: \texttt{a > 0 and b > 0}. Exact-ID overlap is only T1558.003.

\capturebox{2}{Overlap tab in ATT\&CK Navigator and full \texttt{check\_layer.py --overlap} output.}

\capfig{capture2_navigator_overlap.png}

\captionof{figure}{CAPTURE \#2: Meridian and FIN7 (G0046) overlap tab. Exact-ID match T1558.003 Kerberoasting highlighted (score 1, red on default scale).}

\capfig{capture2_overlap.png}

\captionof{figure}{CAPTURE \#2 companion: overlap fraction = 1 of 67 (FIN7 group total).}

\begin{lstlisting}
Layer: 'Meridian and FIN7 (G0046) overlap'
domain: enterprise-attack   ATT&CK v19   Navigator 5.3.2   layer format 4.5

Techniques in BOTH layers (score 1): 1
  T1558.003   credential-access

Techniques scored in either layer: 1
Report the overlap as a fraction of your GROUP layer's total
(the count shown next to X when you selected the group).
\end{lstlisting}

\textbf{Overlap fraction: 1 of 67.}

\textbf{Near-misses} (exact-ID tool miss): Meridian parent T1027 vs FIN7's T1027.010 / T1027.016; Meridian T1071.001 Web Protocols vs FIN7's T1071.004 DNS.

\labstep{Adversary profile}

FIN7 (G0046) is a financially motivated enterprise intrusion group documented across retail, hospitality, and broader industries, with later big-game-hunting activity. Two documented procedures that relate to Meridian's evidence, both recorded on the FIN7 \textbf{group page} (not a campaign page): (1)~T1558.003 Kerberoasting, which is an exact-ID match to Meridian E2; (2)~T1071.004 DNS (Application Layer Protocol), which is the same tactic family as Meridian's scored T1071.001 Web Protocols but a different sub-technique, so it is a near-miss rather than an exact overlap. Meridian's exact-ID overlap with FIN7 is 1 of 67 (only T1558.003). That small fraction shows shared tradecraft, not identity. Meridian's evidence does \textbf{not} support attributing the intrusion to FIN7: there is no FIN7-specific malware family, no confirmed POS/ransomware payload, and no unique infrastructure overlap beyond a documentation-only \texttt{.test} host and RFC~5737 address used in the course lab.

\labstep{Q5. Why might a 119-technique group overlap more without being likelier?}

A group with 119 documented techniques (for example APT29) covers more of the ATT\&CK matrix, so random or common techniques in a small Meridian layer are more likely to collide by chance. Overlap count is not a likelihood score: denser public reporting inflates matches without proving that group operated at Meridian. Attribution still needs unique malware, infrastructure, or procedure-level detail, and an honest statement of what the evidence cannot support.

\labpart{Part C. STIX 2.1 bundle (Tasks 8--10)}

Option~A (file IOC): Indicator pattern uses the E3 SHA-256 of the \texttt{svcupdate.bin} stand-in. Attack-Pattern is T1027 Obfuscated Files or Information (score~$45 \ge 30$; not the lecture's T1558.003). Relationship: \texttt{indicator --indicates--> attack-pattern}.

\begin{lstlisting}
pattern: [file:hashes.'SHA-256' = '3f77970a8c10f1af034e984b3b06f76db1708d377e8c3f9085af95823c9e2d73']
attack-pattern external_id: T1027
url: https://attack.mitre.org/techniques/T1027/
source_name: mitre-attack
\end{lstlisting}

\capturebox{3}{Complete \texttt{check\_stix.py} output ending in RESULT: ALL CHECKS PASSED (\{302\} warning ignored).}

\begin{lstlisting}
Python 3.11.9 | stix2-validator 3.2.0 | stix2 3.0.2
[!] Warning: {302} External reference 'mitre-attack' has a URL but no hash.
[PASS] 1. stix2-validator
[PASS] 2. stix2.parse (strict object model)
[PASS] 3. Lab 5.1 structure and ATT&CK cross-reference

RESULT: ALL CHECKS PASSED
\end{lstlisting}

\capfig{capture3_stix_check.png}

\captionof{figure}{CAPTURE \#3: STIX checker summary figure.}

\labstep{CAPTURE \#4 and Q6: visualization}

\capturebox{4}{OASIS STIX Visualizer graph for \texttt{meridian\_stix.json}: 2 nodes, 1 edge labeled \texttt{indicates}.}

\capfig{capture4_stix_visualizer.png}

\captionof{figure}{CAPTURE \#4: OASIS STIX Visualizer. Indicator \texttt{svcupdate.bin} stand-in SHA-256 indicates Attack-Pattern Obfuscated Files or Information (T1027).}

\textbf{Q6.} The graph shows \textbf{2 nodes} and \textbf{1 edge}. The edge label is \textbf{indicates}. The Relationship object appears as an edge because STIX visualizers render SROs (STIX Relationship Objects) as links between SDOs (STIX Domain Objects). Indicator and Attack-Pattern are domain objects (nodes); the \texttt{indicates} relationship is an SRO, so the tool draws it as the connecting edge rather than as a third node. Do not use the timeline slider when checking properties: it can hide objects by timestamp.

\labpart{Challenge. Opportunistic or targeted?}

\textbf{Position:} Meridian's intrusion looks more like a \emph{targeted} campaign against Meridian than a purely opportunistic smash-and-grab, but the case is not closed.

Two scored techniques support a Meridian-focused operator. T1558.003 (score~95) shows Kerberoasting aimed at \texttt{svc-billingsql}'s MSSQL SPN on the Meridian domain. That is a purposeful credential attack against a named billing/SQL service, not a random internet scan. T1071.001 (score~75) shows long-lived TLS beaconing from BILL-WS-14 to a Meridian-themed relay host (\texttt{sync-update.meridian-relay.test}) with a host-unique JA3. That combination of a Meridian service-account target and Meridian-branded C2 naming is hard to square with an actor who only stumbled onto an open port.

A targeting-relevant Evidence Card fact is E3's embedded strings: \texttt{MeridianLoader/2.1} and \texttt{sync-update.meridian-relay.test} inside the stand-in recovered from BILL-WS-14. That is Meridian-specific branding in the loader/C2 path. From Task~7, the FIN7 overlap is only \textbf{1 of 67} exact technique IDs, which warns against treating a famous group label as proof of identity even when Kerberoasting overlaps.

What would strengthen targeting: packet capture of the TLS sessions showing Meridian data theft, or evidence that \texttt{svc-billingsql} access led to billing records. What would overturn targeting toward opportunistic: discovery that many unrelated orgs hit the same loader hash/JA3 with generic naming, or proof the \texttt{.test} host was only a shared commodity C2 with no Meridian focus. Until then, I read the campaign as Meridian-directed tradecraft with incomplete attribution.

\labpart{References}

\begin{itemize}
  \item MITRE ATT\&CK Enterprise Matrix (v19): \url{https://attack.mitre.org/}
  \item MITRE ATT\&CK Navigator 5.3.2: \url{https://mitre-attack.github.io/attack-navigator/}
  \item OASIS, STIX Version 2.1 (2021)
  \item OASIS STIX Visualizer: \url{https://oasis-open.github.io/cti-stix-visualization/}
  \item Bianco, D.~J. (2013). The Pyramid of Pain.
  \item Althouse, J., Atkinson, J., \& Atkins, J. (2017). JA3: a method for profiling SSL/TLS clients.
\end{itemize}

\end{document}
